% =====================================================================
% tesi/capitoli/31-corpus-e-pianificazione.tex
% =====================================================================
% copre: pokedex-home-completo/LETTURA-DEL-CORPUS.md
% copre: pokedex-home-completo/LETTURA-DEL-CORPUS.md#la-lettura-del-corpus-cluster-per-cluster
% copre: pokedex-home-completo/LETTURA-DEL-CORPUS.md#lo-stato-aggiornato-a-ogni-lotto
% copre: pokedex-home-completo/LETTURA-DEL-CORPUS.md#che-cosa-ha-aggiunto-il-cluster-sull-accesso-alla-banca
% copre: pokedex-home-completo/LETTURA-DEL-CORPUS.md#che-cosa-ha-aggiunto-il-cluster-delle-altre-liste-ed-è-la-misura-di-una-nostra-cecità
% copre: pokedex-home-completo/LETTURA-DEL-CORPUS.md#che-cosa-ha-aggiunto-il-cluster-dei-fiocchi
% copre: pokedex-home-completo/LETTURA-DEL-CORPUS.md#che-cosa-ha-aggiunto-il-cluster-delle-collezioni-di-una-sola-specie
% copre: pokedex-home-completo/LETTURA-DEL-CORPUS.md#una-correzione-che-viene-dall-utente-e-non-dalle-fonti
% copre: pokedex-home-completo/LETTURA-DEL-CORPUS.md#lotto-del-2026-09-09-secondo-i-cataloghi-per-generazione-dalla-prima-alla-terza-e-la-quarta
% copre: pokedex-home-completo/LETTURA-DEL-CORPUS.md#che-cosa-vincola-la-prima-generazione
% copre: pokedex-home-completo/LETTURA-DEL-CORPUS.md#che-cosa-vincola-la-seconda-generazione
% copre: pokedex-home-completo/LETTURA-DEL-CORPUS.md#che-cosa-cambia-nella-terza-generazione-e-riguarda-la-scadenza
% copre: pokedex-home-completo/LETTURA-DEL-CORPUS.md#la-quarta-generazione-e-una-voce-che-si-chiude
% copre: pokedex-home-completo/LETTURA-DEL-CORPUS.md#che-cosa-restava-e-perché-era-rimasto-per-ultimo
% copre: pokedex-home-completo/LETTURA-DEL-CORPUS.md#il-risultato-che-vale-più-dell-intero-lotto-e-che-non-riguarda-i-cromatici
% copre: pokedex-home-completo/LETTURA-DEL-CORPUS.md#il-vincolo-del-trasferitore-che-nessuna-fonte-del-progetto-aveva
% copre: pokedex-home-completo/LETTURA-DEL-CORPUS.md#due-vincoli-di-ottenibilità-che-toccano-il-profilo-di-collezione
% copre: pokedex-home-completo/LETTURA-DEL-CORPUS.md#due-porte-che-si-aprono-nel-deposito-e-che-il-progetto-non-aveva-contato
% copre: pokedex-home-completo/LETTURA-DEL-CORPUS.md#che-cosa-il-resto-del-lotto-ha-confermato-senza-aggiungere
% copre: pokedex-home-completo/LETTURA-DEL-CORPUS.md#lo-stato-del-corpus-dopo-questo-lotto
% copre: pokedex-home-completo/SPOGLIO-CORPUS.md
% copre: pokedex-home-completo/SPOGLIO-CORPUS.md#spoglio-del-corpus-che-cosa-le-fonti-nominano-e-noi-non-abbiamo
% copre: pokedex-home-completo/INDICE-FOGLI-ESTERNI.md
% copre: pokedex-home-completo/INDICE-FOGLI-ESTERNI.md#indice-delle-cartelle-di-calcolo-esterne-della-collezione
% copre: pokedex-home-completo/ROADMAP.md
% copre: pokedex-home-completo/ROADMAP.md#roadmap-cronologica-del-completamento-del-deposito
% copre: pokedex-home-completo/ROADMAP.md#come-leggere-questo-documento
% copre: pokedex-home-completo/ROADMAP.md#il-quadro-dove-siamo-il-2026-09-08
% copre: pokedex-home-completo/ROADMAP.md#fase-a-chiudere-ciò-che-è-già-in-mano-e-che-nessuno-sta-aspettando
% copre: pokedex-home-completo/ROADMAP.md#fase-b-le-due-misure-che-decidono-la-forma-del-lavoro-seguente
% copre: pokedex-home-completo/ROADMAP.md#fase-c-la-produzione-che-resta
% copre: pokedex-home-completo/ROADMAP.md#fase-d-il-trasferimento-che-è-il-collo-di-bottiglia-vero
% copre: pokedex-home-completo/ROADMAP.md#che-cosa-resta-fuori-e-perché
% verificato-al-commit: 80fac1f
% ---------------------------------------------------------------------

\chapter{Leggere un corpus, e pianificare con ciò che se ne ricava}
\label{cap:corpus}

Il capitolo~\ref{cap:fonti} ha stabilito come si valuta una fonte e in quale ordine si crede a due fonti che si contraddicono. Questo capitolo tratta il problema che nasce quando le fonti non sono dieci ma trecento, e la lettura in sequenza cessa di essere possibile: un corpus di duecentonovantadue documenti per sei milioni e duecentomila byte non entra in alcuna sessione di lavoro, e nemmeno in alcuna settimana di lettura umana attenta.

Il corpus di cui si parla discende da un solo messaggio di raccolta, che rinvia a centosettantuno fonti, e dal grafo dei rinvii che parte da quelle, che ne aggiunge millecentoventotto. Il modo in cui è stato attraversato è descritto nel capitolo~\ref{cap:strumenti}; qui interessa il modo in cui è stato \emph{letto}, che è una questione diversa e per la quale questo lavoro ha adottato tre tecniche distinte.

\section{Tre modi di leggere ciò che non si può leggere}

Il primo modo è la lettura per predicato invece che in sequenza. Si costruiscono i vocabolari chiusi del dominio, cioè i nomi delle specie, delle mosse e dei fiocchi letti dalle tabelle del verificatore, e si chiede a un programma che cosa il corpus nomini e i nostri lotti non contengano. È una lettura integrale nel senso che tocca ogni documento, e non lo è nel senso che non ne comprende nessuno: risponde a una domanda sola e la risponde su tutto. I suoi due limiti sono dichiarati e visibili nell'esito, cioè che i nomi corti o ambigui sono esclusi per non produrre falsi positivi e che i nomi scritti in una lingua diversa dall'inglese non sono riconosciuti affatto. La direzione che interessa è una sola, cioè ciò che il corpus nomina e i lotti non hanno; la direzione opposta non è un difetto e non si riferisce, perché il corpus non pretende di essere completo.

Il secondo modo è la lettura per livelli delle cartelle di calcolo, che nel corpus sono la forma in cui la comunità pubblica le proprie enumerazioni. Una cartella comunitaria porta centinaia di migliaia di celle piene e la domanda a cui si risponde aprendola è quasi sempre soltanto quale scheda guardare: si genera dunque uno scheletro di primo livello che per ciascuna scheda dà righe, colonne, celle piene, riempimento e intestazione riconosciuta, e si scende alla lettura vera soltanto sulla scheda che serve. Il riempimento va letto come indizio della forma di una scheda, perché un valore basso indica una tabella sparsa o una scheda di sola prosa e un valore alto una enumerazione densa; le caselle di spunta si contano a parte perché distinguono una enumerazione da leggere dallo stato di avanzamento di chi ha compilato il foglio, che a noi non serve.

Il terzo modo è la lettura vera, fatta a lotti e per cluster, con un documento che si accresce a ogni lotto e che dice per ciascun cluster che cosa ha aggiunto, che cosa ha corretto e che cosa resta non letto con il motivo. Non è un censimento generato ed è la parte che nessun programma può fare, perché consiste nel riconoscere che una riga di una fonte tocca una nostra affermazione. L'ordine dei cluster non è quello del messaggio di raccolta ma quello dell'utilità per il lavoro in corso, che nella fase attuale coincide con l'ordine della scadenza.

\section{Che cosa ha prodotto la lettura, cluster per cluster}

Il cluster sull'accesso alla banca ha portato la voce più importante di tutta la lettura, cioè la regola di ammissione del deposito, che il capitolo~\ref{cap:catena} tratta per esteso perché cambia la pianificazione. Le altre tre voci hanno reso meno di quanto il titolo promettesse, e vale registrarlo perché un cluster povero è un dato: la guida al modding della console è quella che il progetto già segue, una tabella di connettività sta su un servizio di condivisione che non la lascia scaricare, e la guida all'allestimento serve a chi deve comprare l'hardware e non a chi lo ha già. Una cosa quella guida però la dice e il lavoro non l'aveva scritta in forma esplicita: la catena richiede due console fisiche nello stesso momento per il passaggio dalla quarta alla quinta generazione, perché quel trasferimento avviene fra due apparecchi in comunicazione e non dentro una console sola. Era un vincolo implicitamente soddisfatto e mai enunciato.

Il cluster delle altre liste ha prodotto la misura di una nostra cecità, ed è quella descritta nel capitolo~\ref{cap:assi} sugli identificativi di allenatore documentati. Vale qui la nota di metodo che l'accompagna: quella fonte non era stata scaricata dall'attraversamento automatico e una richiesta locale l'ha restituita al primo tentativo, il che è la conferma pratica della prescrizione secondo cui la via più economica va tentata quando si cataloga e non un giorno dopo. Due voci dello stesso cluster non si recuperano e sono catalogate con il motivo, perché stanno su un servizio di condivisione che rifiuta sia l'esportazione sia lo scaricamento diretto: entrambe sono materiale che una persona scarica dal browser in un minuto, e stanno fra le cose attese con la domanda a cui servono.

Il cluster dei fiocchi ha portato poco in contenuto e una cosa in metodo. Il manuale comunitario è ospitato su un servizio che rende le pagine nel browser, quindi l'attraversamento ne ha salvato il solo guscio, cioè mille e duecento byte di intestazione e nessun testo. Ciò che quel guscio porta è però significativo, perché è l'unica riga che il servizio scrive nel documento iniziale ed è un avviso: la banca chiude alla data nota, chi vuole completare quella generazione lo faccia prima, ed esiste una pagina di preparazione d'emergenza per chi voglia chiudere l'obiettivo entro la scadenza. Ne segue che quel manuale è una fonte da recuperare e non da archiviare, perché ha una pagina dedicata esattamente al nostro problema.

Il cluster delle collezioni di una sola specie documenta gli assi che questo lavoro chiama minori e che una decisione dedicata prescrive di misurare invece di moltiplicare. La voce che serve è quella sui livelli minimi per sfera, che per ciascuna specie dichiara il livello più basso a cui la si può avere in ciascuna sfera distinguendo l'esemplare cresciuto da quello selvatico: è esattamente la tavola che serve a comporre una coppia fra specie e sfera che il verificatore accetti. Anche qui va dichiarato un limite nostro, perché l'estrattore ha perso le intestazioni delle colonne conservando i numeri, e la tavola è oggi illeggibile nel derivato mentre il grezzo, conservato accanto, la contiene ancora.

\section{I cataloghi per generazione, letti nell'ordine della scadenza}

Il lotto di lettura successivo ha seguito l'ordine imposto dalla scadenza, cioè le generazioni la cui via verso il deposito passa dalla banca. La fonte portante sono tre guide al catalogo per regione di origine, aggiornate rispettivamente al novembre 2025 le prime due e all'agosto 2026 la terza, cioè dopo l'annuncio della chiusura.

Sulla prima generazione \cite{guide-origine-gen1} il fatto che tocca il lavoro più di ogni altro riguarda la cromaticità: un esemplare selvatico non può essere legittimamente cromatico, mentre lo possono essere i doni, gli scambi in gioco e gli incontri fissi. È un vincolo di legalità e non una probabilità, quindi ogni voce cromatica composta per quella generazione deve appartenere a una di quelle tre classi, e le sedici voci cromatiche del nostro lotto vanno riviste con questo criterio. La stessa guida conferma per via indipendente la verifica sui premi dei tornei chiusa poche ore prima e la rende più precisa, aggiungendo che per iniezione passano tutti tranne un esemplare che conosce una certa mossa di supporto, perché il verificatore del trasferimento non la ammette su quella specie: la voce corrispondente del nostro lotto va marcata come non trasferibile nemmeno per iniezione. Sugli esclusivi di versione la guida enumera ciò che il lavoro non aveva scritto, e ne discende che il catalogo di quella generazione richiede tre versioni e non una, cosa che il progetto sapeva per la terza generazione e non aveva enunciato per la prima.

Sulla seconda generazione \cite{guide-origine-gen2} il limite da conoscere prima di comporre è che alcune classi di specie non si ottengono legittimamente in quei titoli, cioè gli iniziali della prima generazione, i fossili tranne uno, i tre uccelli leggendari e i due mitici: chi le volesse con provenienza di quella regione deve passare per la generazione precedente. Sull'asse delle forme arriva un dato preciso che tocca una sfida del deposito, cioè che la specie dalle molte forme ne ha ventisei in quella generazione e che le due di punteggiatura compaiono per la prima volta nelle riedizioni della terza: la sfida che ne chiede ventotto non è dunque soddisfacibile con la sola seconda generazione, e questo spiega perché la nostra enumerazione ne trovasse ventotto dove la generazione di origine ne ha ventisei.

Sulla terza generazione \cite{guide-origine-gen3} la guida porta il fatto che sposta un risultato misurato: la versione per console corrente dei due titoli di Kanto include l'evento dell'isola della nascita, quindi la specie che ne viene si ottiene là, si può cercare cromatica e registra Kanto come regione. Quella specie era una delle quattro che una misura precedente aveva trovato prive di qualunque incontro nei giochi moderni: se la fonte è esatta, quella riga va corretta e una specie in meno dipende dalla catena storica. Il resto della guida completa l'asse delle uova con quattro voci di cui il lavoro aveva registrato l'esistenza e non il dettaglio tecnico, cioè che ciascuna porta una mossa esclusiva e che soltanto una deve dimenticarla per essere trasferita: le altre tre la conservano fino al deposito, e sono quindi collezionabili con una proprietà che nessun'altra via riproduce.

Sulla quarta generazione si chiude un punto che era in sospeso, ed è il caso in cui due vie studiate separatamente si toccano \cite{sinjoh-ruins}. Il terzetto del tempo e dello spazio al livello uno si ottiene portando un esemplare mitico di quarta generazione in una riedizione della seconda, tenendolo come solo esemplare in squadra, entrando in un luogo dedicato e scegliendo uno dei tre cerchi. Ne segue che la voce non è sbloccata da un oggetto ma dal possesso di quell'esemplare, e che quell'esemplare a sua volta viene da un oggetto mai distribuito ufficialmente oppure dalla distribuzione di un rivenditore, che il servizio in rete ricostruito del capitolo~\ref{cap:dentro} rende di nuovo ottenibile.

Va infine registrata una correzione che non viene dalle fonti ma dall'utente, perché il modo in cui è arrivata è istruttivo. Il vincolo di calendario sollevato sul piano a pagamento era già soddisfatto, perché il piano è attivo e il servizio di deposito intermedio è gratuito da anni; il contatore di giorni che quell'applicazione mostra è il residuo di quella gratuità e non un abbonamento da rinnovare, ed era già stato registrato come artefatto privo di significato. Resta della questione la sola parte che conta, cioè la capienza, che va contata nel piano perché il perimetro aperto dalle decisioni recenti è più grande di essa.

\section{La pianificazione, e perché l'ordine non è quello dell'importanza}

Il prodotto di tutta questa lettura non è un elenco di cose aperte, che vive altrove, ma una sequenza. La differenza fra le due è quella che rende una pianificazione utilizzabile: un elenco dice che cosa manca, una sequenza dice che cosa fare adesso e perché non altro.

Ogni passo dichiara tre cose che di solito restano implicite e che, restando implicite, fanno perdere tempo. Chi lo esegue, perché alcuni passi sono lavoro di sessione e altri richiedono una persona davanti a una console. Che cosa lo sblocca, perché un passo eseguito prima della sua dipendenza va rifatto. E che cosa il suo esito cambia, perché un passo il cui esito non cambia nulla è un passo che si può rimandare per sempre. L'ordine non è quello dell'importanza ma quello del rendimento: un passo che costa poco e che, comunque vada, elimina un'incertezza su cui altri passi poggiano viene prima di un passo importante e isolato.

Il quadro da cui la sequenza parte è misurato e non stimato. Cinque lotti sono prodotti con stato di giudizio diverso: centosessantacinque esemplari per le prime due generazioni con tre difetti corretti dopo il giudizio esterno, centosettantadue voci tutte conformi per la terza, duecentoquarantasette per la quarta di cui duecentodiciannove conformi e ventotto rifiutati per la barriera di lingua, settecento tutti conformi per la quinta, e dieci incontri sbloccati da un oggetto in attesa di giudizio. Due lotti sono progettati e non prodotti, cioè le distribuzioni che non lasciano una carta e le quattordici specie che soltanto un archivio esterno conosce. E il collo di bottiglia non è più la produzione ma il trasferimento.

La prima fase chiude ciò che è già in mano e che nessuno sta aspettando, cioè il giudizio esterno sul lotto degli incontri, la lettura dei cluster prioritari del corpus e il completamento del censimento delle fonti. Sono passi che non dipendono da nulla, e la ragione per cui vanno per primi è che ogni giorno in cui restano aperti è un giorno in cui qualche altro passo poggia su un'incertezza evitabile.

La seconda fase sono due misure che decidono la forma del lavoro seguente e costano poco: che cosa il canale ricostruito distribuisca oggi, e se quel canale attraversi la barriera di lingua della quarta generazione. La prima decide se una parte della produzione rimanente abbia ancora senso; la seconda decide se un sottoproblema si risolva con un servizio già in piedi invece che con una stazione in emulazione, il che ne cambia il costo di un ordine di grandezza.

La terza fase è la produzione che resta, e comprende il lotto delle distribuzioni senza carta, il confronto fra un esemplare ricevuto e il suo gemello composto, e le specie che soltanto l'archivio esterno conosce. Il passo di maggior valore conoscitivo dell'intera sequenza è il secondo, e la ragione è che un esemplare ricevuto davvero dal gioco è verità di riferimento: confrontarlo campo per campo con il nostro gemello dice se il modello sia giusto, cosa che nessun verificatore può dire, perché un verificatore controlla la conformità a un modello e non il modello.

La quarta fase è il trasferimento, che è il collo di bottiglia vero, e i suoi quattro passi sono la conferma del catalogo nazionale sui salvataggi di destinazione, il cronometraggio di una sessione del parco di migrazione, l'allestimento della via in emulazione e il trasferimento del lotto di quinta generazione, che va per primo perché è il più grosso e il più economico.

Resta da dichiarare che cosa la sequenza lasci fuori, perché una pianificazione che non lo dica sembra completa e non lo è. Gli esemplari che una macchina nascosta blocca non vi entrano, e non è una dimenticanza: si conservano come file e si attende una catena che quel vincolo non abbia, quindi la loro coda si riesamina quando una catena nuova diventa disponibile e non a una data. La forma dalle orecchie a punta non vi entra perché la questione è chiusa in negativo, dato che il trasferimento la rifiuta esplicitamente. E le forme di sola battaglia e quelle totemiche non vi entrano perché sono irraggiungibili sempre e non per scadenza, che è una distinzione che questo lavoro ha dovuto correggere una volta e che non conviene perdere.

\section{Lo stato di lettura come dato invece che come prosa}

La tabella con cui il registro di lettura dichiara a che punto sia ciascun cluster era scritta per una persona, e i nomi dei cluster vi comparivano abbreviati o raggruppati in una frase. Dal momento in cui una seconda vista ha avuto bisogno di quello stato, cioè l'indice unico dentro il registro delle fonti, quella libertà è diventata un difetto: un nome abbreviato non corrisponde ad alcun cluster del censimento, e la vista mostrava senza stato quattro quinti dei cluster già letti.

La correzione è minima e vale come principio generale su questo genere di documenti. Il nome del cluster in quella tabella è ora esattamente quello del censimento, e quando una riga vale per più cluster i nomi si elencano separati da virgola; la riga di prosa che diceva quanti cluster restassero è stata tolta, perché quel numero lo calcola la vista e una riga di prosa che porta un conto invecchia in silenzio. Il documento resta leggibile da una persona e diventa leggibile da un programma senza che si sia introdotto un formato nuovo: è la stessa scelta già fatta per le dichiarazioni di copertura in testa ai capitoli, cioè un commento che una persona legge e uno strumento verifica.

\section{Un corpus che si credeva letto, e la macchina che lo rendeva illeggibile}

La lettura di un corpus poggia su un presupposto che nessuno enuncia, cioè che l'elenco di ciò che non si è potuto leggere dica la verità sul perché. Quando quell'elenco sbaglia la ragione, il difetto non si manifesta come un errore ma come una assenza che sembra motivata, e questo è il modo in cui un corpus perde un terzo di sé senza che nulla protesti.

Il lotto di lettura dichiarato contava ventiquattro voci e ne risultavano scaricate sette; per quattordici delle diciassette mancanti il censimento riportava che il file delle regole di accesso del loro sito non era leggibile per indisponibilità del servizio. Tre verifiche in sequenza hanno smentito la lettura ovvia di quella riga. Il file delle regole risponde oggi e non vieta le pagine escluse, anzi le consente con una direttiva che le nomina esplicitamente. La medesima ragione aveva fermato centotrenta nodi su ventotto siti in una sola corsa, fra cui un'enciclopedia generalista e un archivio tecnico, e ventotto guasti simultanei di servizi indipendenti sono meno probabili di un guasto solo dal lato di chi chiede. E la richiesta che falliva non falliva per un codice di risposta, ma prima, nella verifica del certificato.

La causa è locale e misurabile. L'archivio delle radici di fiducia di questa macchina contiene nove certificati scaduti, e fra essi la copia di una radice intermedia scaduta nel settembre 2025. I siti che si appoggiano a quell'autorità servono una catena che termina in quella radice controfirmata da una radice superiore valida per un altro decennio; la verifica, trovando la prima fra quelle di cui si fida, vi si ferma e la dichiara scaduta invece di proseguire lungo il ramo controfirmato, che è valido. Un browser e un recupero da riga di comando non se ne accorgono, perché usano un archivio diverso; un programma che usa l'archivio di sistema sì. Il messaggio che ne risultava nominava il sito remoto mentre il difetto stava nella macchina locale, e questa è la parte del difetto che è costata di più, perché una diagnosi sbagliata chiude la ricerca invece di aprirla.

Ne discende una regola di metodo che generalizza quella già adottata per le fonti non recuperabili. Quando più tentativi falliscono con il medesimo esito, prima di cercarne uno in più conviene domandarsi che cosa abbiano in comune: se condividono l'estremo che rifiuta, non sono tentativi diversi ma lo stesso tentativo ripetuto. Nel caso delle fonti quell'estremo era il sito; qui è la macchina che chiede, e il criterio è lo stesso applicato un gradino più in basso.

La riparazione ha due parti indipendenti, e la seconda non sarebbe bastata senza la prima. Il rifiuto dovuto a un file di regole illeggibile è ora distinto dal divieto scritto: il primo è una proprietà dell'istante e il nodo va fra quelli in attesa, dove il programma metteva già ciò che un tetto aveva escluso e per la medesima ragione, mentre il secondo è una proprietà del sito e la registrazione definitiva è corretta. La verifica dei certificati usa inoltre un archivio di radici corrente invece di quello del sistema, e quando fallisce lo dichiara con quelle parole anziché chiamarlo guasto di rete. Ciò che non è stato fatto va detto, perché era la scorciatoia disponibile: la verifica non è stata disattivata e non esiste un'opzione per disattivarla, dato che il difetto stava nell'elenco di chi ci si fida e non nel fatto di fidarsi.

Il recupero delle corse già concluse è stato reso un'operazione dello strumento anziché un intervento a mano sul suo stato, perché il caso tornerà; ed è separabile dal proseguimento della frontiera, perché riparare un difetto e continuare una lettura sono due cose diverse che è utile poter fare una alla volta. L'esito è che il corpus passa da duecentonovantadue a quattrocentocinque documenti, con settantaquattro pagine enciclopediche recuperate, e che il lotto dichiarato passa da sette voci leggibili a quindici.

\section{Che cosa i cataloghi della coda hanno aggiunto}

Le voci recuperate hanno reso leggibile la coda dei cataloghi per generazione, e il rendimento è stato disuguale in un modo che vale registrare, perché è la ragione per cui l'ordine di lettura non coincide con l'ordine del post.

La quinta generazione ha trasformato in elenco una voce di lavoro che il progetto portava come conteggio. Il dispositivo che estrae esemplari da una dimensione onirica ne consegna ventisei specie in quattro insiemi che si sbloccano in sequenza, e quattro proprietà li rendono un collezionabile distinto invece che una scorciatoia per specie già coperte: portano la propria abilità nascosta con una sola eccezione, vengono in una sfera che appartiene a quella generazione, hanno come luogo d'incontro il dispositivo stesso, e hanno un livello che non è fissato dalla tavola ma dal numero di medaglie possedute nel gioco che li riceve \cite{bulbapedia-dream-radar}. L'ultima proprietà corregge un'aspettativa: il livello cinque, che questo lavoro aveva indicato come il pregio di quella via, non è una proprietà dell'esemplare ma una conseguenza del momento in cui lo si riceve, e va quindi ottenuto su un salvataggio senza medaglie. La stessa pagina registra inoltre che le schede di gioco della generazione precedente, inserite per sbloccare i cinque leggendari, funzionano indipendentemente dalla propria lingua: è un precedente e non una soluzione per la barriera di lingua che blocca ventotto voci altrove, ma è la prima volta che si incontra un meccanismo di quella generazione dichiaratamente indifferente alla lingua di ciò che legge.

Gli esemplari appartenenti al personaggio antagonista della medesima generazione sono quindici, contati sulla pagina, il che conferma per via indipendente la misura condotta sulle tavole del verificatore e che aveva corretto un trentasei portato in memoria \cite{bulbapedia-n-pokemon}. La fonte aggiunge i tratti che li rendono irriproducibili, cioè un identificativo fisso con un allenatore proprio, trenta in ogni valore individuale, natura fissata, impossibilità di ricevere un soprannome e cromaticità bloccata nel codice; e ne toglie uno che si sarebbe potuto assumere, perché la sfera non è fissata ma è quella che il giocatore usa.

La settima generazione ha prodotto una conferma che chiude una questione invece di aprirne una. Gli esemplari di taglia totemica si ottengono soltanto nelle riedizioni, in cambio di adesivi, e se trasferiti nel deposito vengono riportati alla taglia ordinaria \cite{bulbapedia-totem}: non sono dunque irrecuperabili dopo la scadenza ma irraggiungibili sempre, che è esattamente la correzione registrata da questo lavoro su un'altra base. Una fonte che non conosce il ragionamento e vi arriva per un'altra strada è ciò che serve a promuovere una conclusione a fatto. La stessa generazione ha invece prodotto un insieme che il progetto non aveva, cioè quarantasette specie non native raggiungibili una per giorno e per isola dopo aver accumulato cento punti, ciascuna con una mossa che a quel livello non avrebbe e che ne costituisce la traccia \cite{bulbapedia-island-scan}; l'insieme contiene i capostipiti di ogni regione precedente, il che ne fa una via alla loro cattura dentro quella generazione anziché un trasferimento.

\section{Un lotto che non è un cluster, e perché vale di più}

I lotti di lettura descritti finora erano cluster del post di raccolta, affrontati in ordine di utilità. Il lotto che segue il recupero non lo è, e la differenza non è organizzativa ma conoscitiva: fra le pagine rientrate ce ne sono nove che il lavoro attendeva da giorni e che stanno a profondità due nel grafo dei rinvii anziché dentro un cluster. Il fatto è esso stesso una misura sul corpus, e dice che la parte utile non coincide con la parte citata: un post di raccolta indica dove guardare, non che cosa serve.

L'esito di quel lotto sul primo anello della catena è descritto nel capitolo precedente. Qui interessano le altre tre cose che ha prodotto, perché ciascuna corregge o completa un lavoro che era fermo.

La prima è una correzione a una procedura registrata due settimane prima sulla base di un cluster. L'evento che consegna i tre draghi al livello uno non ne consegna tre ma uno, scelto fra i tre; si può attivare due volte in tutto, una con un esemplare da distribuzione e una con un esemplare proveniente da un luogo raggiungibile solo con un oggetto mai distribuito, dopo di che diventa irraggiungibile per sempre \cite{bulbapedia-sinjoh}. Poiché quell'oggetto non fu mai rilasciato e il relativo esemplare è già registrato come non producibile in forma conforme, su un salvataggio si ottiene un drago solo, e ottenerne tre richiede tre salvataggi. La corrispondenza fra i tre cerchi e le tre specie, che la nota precedente riportava, è invece confermata: l'errore non stava nella mappatura ma nel numero, che la nota lasciava implicito.

La seconda riguarda la classe degli scambi in gioco, enumerata su fonte di primo livello ma non ancora producibile. La fonte enciclopedica dice due cose che le tabelle del verificatore non contengono nella forma necessaria \cite{bulbapedia-in-game-trade}. Per la prima generazione, identificativi e valori individuali sono casuali e l'allenatore è una stringa di controllo sensibile alla lingua: la fedeltà è dunque banale nello stesso senso già stabilito per le distribuzioni della medesima generazione, e ne discende che quegli esemplari possono risultare cromatici una volta trasferiti. Il sistema di trasferimento converte però quella stringa in modo permanente in una stringa letterale nella lingua del gioco di provenienza, il che è una trasformazione deterministica che un generatore deve riprodurre per ottenere la forma corretta dopo il passaggio. Resta una lacuna che tocca la producibilità e che va dichiarata: un esemplare di questa classe è definito dal nome dell'allenatore e dal soprannome, entrambi specifici della lingua, mentre il censimento prodotto da questo lavoro porta le voci e non le stringhe.

Quella lacuna si è chiusa come fonte nella medesima giornata, con un recupero mirato della pagina che elenca quei nomi nelle lingue diverse dall'inglese e dal giapponese \cite{bulbapedia-scambi-lingue}. Vi compaiono, per ciascuno scambio, il soprannome in sei lingue e separatamente il nome dell'allenatore in sette, e la colonna che serve a questo lavoro è quella italiana, poiché le cartucce su cui la collezione si costruisce sono italiane. La distinzione fra fonte e lavoro va però tenuta: il dato è disponibile, il censimento non lo incorpora ancora, e l'incorporazione è deterministica e appartiene al programma che genera quel censimento, non a una lettura.

La terza è un criterio, e riclassifica otto delle undici classi che un capitolo precedente aveva raccolto sotto l'etichetta della via chiusa \cite{bulbapedia-vie-chiuse}. Una sola è chiusa dalla cessazione di un servizio: era un titolo scaricabile e non un supporto fisico, e non è più acquistabile dalla chiusura del negozio della console nel gennaio 2019. Le altre sette dipendono da un supporto fisico e non da una rete, compresa quella che il nome fa sembrare la più legata a un servizio, dove nelle edizioni localizzate le missioni sono incluse nella cartuccia e si sbloccano con una parola d'ordine resa pubblica all'epoca. Il criterio che ne discende vale oltre le otto: una via non è chiusa perché è vecchia, ma perché il servizio che la reggeva è spento, e in quel caso l'esclusione è irreversibile e giustifica la produzione, oppure perché il supporto che la reggeva è raro, e in quel caso la questione è di perimetro e di spesa e la decisione non è tecnica. Ciò che queste pagine aggiungono non è l'enumerazione, che il lavoro possedeva già, ma la condizione di ciascuna, che è precisamente ciò che il controllo all'indietro richiedeva.

Va infine registrato che l'indice delle distribuzioni di evento, una delle due voci che un cluster dichiarava non recuperabili, è un indice e non una lista: enumera sottoliste separate per generazione, per modo di consegna e per lingua, e le lingue comprendono quella delle cartucce possedute \cite{bulbapedia-indice-distribuzioni}. Non cambia i byte di alcun esemplare e quindi non tocca la conformità; cambia che cosa si possa affermare sulla legittimità di una collezione costruita su cartucce di una lingua determinata, perché nessuna delle due fonti dell'asse degli eventi dichiara in quali paesi una distribuzione sia stata effettivamente tenuta.

\section{Un difetto che reggeva l'economia della raccolta}

Non tutti i difetti vanno corretti appena trovati, e questo capitolo deve un caso in cui la correzione, presa da sola, avrebbe peggiorato il risultato.

Scrivendo il collaudo di una funzione nuova del lettore del corpus, una prova pretendeva che la mappa registrasse il rinvio da una pagina a ciò che essa cita, ed è fallita. La causa era nell'estrazione dei collegamenti, che riconosceva tre forme sintattiche e non la quarta: quella prodotta dall'estrattore da HTML del medesimo programma, cioè il testo dell'ancora seguito dall'indirizzo fra parentesi tonde. La misura su una pagina enciclopedica reale quantifica il difetto senza ambiguità, con sessanta collegamenti presenti in quella forma e due estratti, entrambi provenienti dall'intestazione che lo strumento scrive di propria mano. Ne segue che la raccolta non si era quasi mai espansa attraverso una pagina esterna ma quasi soltanto attraverso i post, il che spiega retrospettivamente perché le pagine più utili si trovassero tutte alla medesima distanza dal punto di partenza ed elencate da un post, anziché raggiunte l'una dall'altra.

La correzione sta in una riga, e la tentazione è rilanciare la raccolta per recuperare ciò che mancava. La misura dice di non farlo, ed è una misura che non costa alcuna richiesta di rete perché si conduce sui soli documenti già presenti: le pagine esterne già raccolte rinviano a oltre quattromila nodi mai visti, di cui più di tremila verrebbero effettivamente scaricati. Guardare la loro provenienza decide la questione, perché più di mille puntano all'archivio di immagini del medesimo sito, quasi mille a un'enciclopedia generalista, centinaia a social e a negozi di applicazioni. Non è materiale di ricerca: è la navigazione di una wiki, che per costruzione rinvia a tutta sé stessa.

Il difetto, in altre parole, era portante. L'economia della raccolta dipendeva da esso, e la proprietà che rendeva il corpus maneggevole, cioè che una pagina esterna fosse una foglia, non era una scelta di progetto ma un effetto collaterale di un errore. La soluzione adottata separa le due cose: le pagine esterne restano foglie, ma per decisione dichiarata e reversibile con un'opzione, mentre i rinvii vengono comunque registrati nella mappa, perché una mappa deve dire che cosa un documento cita anche quando si è deciso di non seguirlo. Ciò che servirebbe per espandere davvero non è l'opzione ma un criterio che separi un rinvio di contenuto da uno di navigazione, e quel criterio non esiste ancora.

La lezione generalizza oltre il caso, e riguarda il modo in cui si legge un difetto in un sistema che funziona. Quando un comportamento desiderabile discende da un errore, correggere l'errore rimuove il comportamento; la domanda da porsi prima di correggere non è se il codice sia giusto, ma che cosa smetterà di essere vero quando lo sarà.

\section{Una fonte di prima parte che cambia un denominatore}

L'asse dei marchi era stato misurato su un archivio di terze parti, che è la migliore fonte disponibile finché non ne compare una di prima parte. Ne è comparsa una sotto forma di avviso mostrato dal servizio di deposito stesso, consegnato come immagine, e dice tre cose che l'archivio non diceva \cite{avviso-home-mightiest}.

La prima è che le edizioni di questi incontri si ripetono: l'archivio registra un'edizione della medesima specie diciotto mesi prima, e l'avviso ne annuncia il ritorno. La seconda, che è quella che cambia il conto, è che l'esemplare si cattura una sola volta per salvataggio e che chi lo abbia già ottenuto in un'edizione precedente non lo riottiene. Ne segue che l'unità collezionabile è la specie che porta il marchio e non l'edizione che lo distribuisce, e che contare le edizioni produrrebbe un totale più grande del vero: il lavoro di enumerazione lasciato aperto dal capitolo precedente va dunque impostato sul denominatore giusto prima di essere svolto, non dopo. La terza è che l'editore stesso dichiara che quell'esemplare potrebbe comparire in eventi futuri o diventare ottenibile per altre vie, il che indebolisce l'equazione fra edizione chiusa e via chiusa su cui poggia parte del criterio di produzione.

Vale infine registrare una proprietà della fonte derivata, perché è generale e non aneddotica: la copia dell'archivio conservata da questo lavoro è di pochi giorni anteriore all'avviso e si ferma all'edizione precedente. Non è un difetto della copia ma la sua età, e ne discende che il totale letto là è un conteggio a una data e non un numero. Un'enumerazione che cresce va riletta con una cadenza dichiarata, e dichiararla è preferibile a scoprire per caso che è invecchiata.

\section{Uno strumento rotto dalla nascita, e una scappatoia corroborata due volte}

Il lotto di lettura successivo si apriva su un comando di ripresa dell'attraversamento del corpus mai invocato con successo prima d'ora, e la prima cosa che ha prodotto non è stata una lettura ma la scoperta che il comando stesso era guasto: un parametro dichiarato nel sottocomando e letto da una parte del programma che non lo riceveva, un errore immediato a ogni invocazione, presente da quando quel comando esiste. La correzione sta in una riga; il collaudo completo, rilanciato per intero e non solo sulla parte toccata, conferma che nessun altro presidio ne dipendeva in modo fragile.

Con lo strumento funzionante, la stima di quanto restasse da leggere si è rivelata anch'essa sbagliata, e nel verso che sottraeva lavoro anziché aggiungerlo: i nodi effettivamente pendenti erano più del previsto, e leggerli ne ha scoperti altri ancora, portando il numero di post scaricati quasi al triplo di quelli letti fino a quel momento. Tale crescita non è un difetto della stima precedente ma una proprietà del grafo dei rinvii quando lo si attraversa per la prima volta con lo strumento corretto: ogni nodo nuovo può rinviare a nodi che il tetto precedente non aveva mai visto.

La lettura mirata su questo lotto ha portato una seconda fonte indipendente sulla scappatoia del tracciatore discussa nel capitolo~\ref{cap:catena}, con quattro testimonianze concordi e una spiegazione tecnica più precisa della prima, cioè che il titolo in questione condivide con la banca e con il gioco per dispositivo mobile la medesima struttura dati di origine invece di averne una propria, il che la rende un probabile effetto collaterale tecnico e non una scelta deliberata. La stessa fonte porta anche un aneddoto isolato e messo in dubbio nel proprio stesso filo di discussione, che se fosse vero contraddirebbe uno dei limiti già registrati: non è stato promosso a fatto, secondo il medesimo criterio con cui questo lavoro tratta ogni testimonianza singola non corroborata.

Una terza corroborazione, infine, non veniva dal corpus ma da un'altra parte di questo stesso lavoro: la scheda che studia i servizi di scambio automatico su console moderna, scritta mesi prima studiando un canale differente, descrive il medesimo meccanismo del tracciatore in modo indipendente, senza nominare la scappatoia specifica ma confermandone i tratti strutturali. Tre fonti indipendenti, di autori e momenti diversi, che convergono sulla stessa meccanica sono il grado di corroborazione più alto che questo lavoro abbia raggiunto su un dettaglio tecnico non documentato dal titolare del servizio; non bastano tuttavia a rendere praticabile la scappatoia, la cui chiusura in negativo per ragioni pratiche indipendenti è discussa nel capitolo~\ref{cap:catena}.

Un quarto giro di lettura, sui due candidati di punteggio più alto rimasti dal giro precedente, ha aggiunto un vincolo tecnico invece di un'ulteriore corroborazione. Un commento distingue due tecniche di famiglia diversa che è facile confondere: l'iniezione di un dono misterioso tramite l'editor del verificatore replica davvero la distribuzione originale scrivendo il dato nel formato che il gioco si aspetterebbe da un servizio vero, mentre l'esecuzione di codice arbitrario, già nota a questo lavoro, si limita a far comparire l'oggetto dell'evento senza replicare il meccanismo della sua consegna: due vie che possono produrre lo stesso esemplare finale per strade di legittimità diverse. Un secondo commento, indipendente dal primo, conferma un vincolo che nessuna fonte precedente aveva reso esplicito: estrarre un file di salvataggio da una console portatile moderna richiede una console modificata. Non conferma né smentisce la scappatoia del tracciatore, ma stabilisce che farne uso per un esemplare di questo lavoro comporta la stessa esposizione dell'account già registrata come decisione di perimetro aperta altrove in questo lavoro, e non è dunque un'operazione a costo zero anche ammettendo che la scappatoia regga.

\section{Un catalogo che si credeva incompleto, e non lo era: era il generatore a non implementarlo}

Una domanda posta a valle della produzione, non della lettura, ha riportato l'attenzione su un difetto che il corpus non poteva rivelare perché non lo riguardava: quattro distribuzioni note nel catalogo generato di terza generazione, tutte uova consegnate a un allenatore giapponese secondo un metodo di generazione particolare, mancavano dal lotto composto. Non erano assenti dal catalogo, come una prima verifica frettolosa avrebbe potuto concludere, ed è la ragione per cui questo capitolo la registra: erano assenti dal lotto perché il generatore riconosceva una variante del metodo e non l'altra, senza dichiararlo come esclusione nominata nel modo in cui dichiarava l'unica altra voce non implementata. Il conto tornava comunque, perché cinque voci mancanti su centosettantasette erano esattamente il canale non implementato più queste quattro, ma un conto che torna per la ragione sbagliata è un conto che nasconde un difetto invece di confermarne l'assenza. La correzione, e il lotto rigenerato che ne risulta, sono descritte nel capitolo~\ref{cap:distribuzioni}.

Un secondo effetto di questo stesso episodio ha riaperto una domanda che questo capitolo aveva lasciato scritta come da verificare: se un esemplare non ancora schiuso attraversi il primo anello della catena di trasferimento. Le due pagine di riferimento sul tema, già nel corpus, davano una risposta parziale e asimmetrica: quella sul secondo anello esclude le uova esplicitamente, quella sul primo descrive come si stabilisca il gioco di origine di un esemplare nato da uovo senza mai dichiarare le uova escluse dal passaggio stesso. La domanda restava aperta perché una lettura più attenta della prosa non poteva chiuderla, e si chiude infatti non nel corpus ma con una prova diretta sul verificatore, descritta nel capitolo~\ref{cap:catena}: un esemplare non schiuso attraversa il primo anello, e ne esce schiuso, coerentemente con un gioco vero dove un uovo non può attraversarlo restando tale.

\section{La chiusura del corpus, e ciò che la coda conteneva}

La lettura del corpus si è chiusa il 2026-09-16 sui diciassette cluster rimasti, sessanta voci in tutto, delle quali trentaquattro erano scaricate e leggibili, ventiquattro catalogate con il proprio motivo, quasi tutte perché in forma di video, e due fuori dal grafo della corsa. Erano rimasti per ultimi con una motivazione dichiarata e in sé corretta, cioè che riguardano quasi tutti la caccia alle forme cromatiche e servono soltanto se il profilo di collezione scelto le comprende. Quella motivazione è però diventata un ostacolo nel momento in cui la decisione sul profilo ha continuato a restare aperta: un corpus letto a metà costringe a tornare sulle medesime domande, e il costo cumulato di quei ritorni ha superato il costo della lettura. È un criterio di priorità che vale registrare perché la sua forma si ripete, cioè rimandare ciò che dipende da una decisione che non arriva finisce per costare più del farlo.

Il risultato più importante del lotto non riguarda affatto i cromatici. La pagina enciclopedica delle forme cromatiche non ottenibili porta, in una nota sul mitico della foresta, la conferma indipendente di ciò che la medesima giornata aveva stabilito leggendo il sorgente del verificatore: gli esemplari delle riedizioni per console corrente possono salire soltanto alla settima generazione o oltre passando per la banca, mentre quelli ottenuti sulle cartucce fisiche delle prime due generazioni non possono salire alla terza né a una successiva. È la terza prova indipendente dello stesso fatto, dopo la condizione esplicita nella funzione che decide l'ammissibilità di una conversione e la tabella delle quattro catene, e chiude la questione aperta nel capitolo~\ref{cap:catena}.

Il manuale della caccia per le prime due generazioni aggiunge il secondo tempo di quella conferma, e lo fa dal lato della via che invece esiste, documentando per esteso un dettaglio che il sorgente esprime come una chiamata di funzione: la natura di un esemplare trasferito si ricava dai punti esperienza come resto della divisione per venticinque. Ne discende una conseguenza operativa che trasforma un effetto collaterale in uno strumento, perché la natura di ciò che attraversa quell'anello non è casuale ma interamente determinata da un campo che questo lavoro compone, e dunque è scegliibile.

Sulla stessa via compare un vincolo puntuale che nessuna fonte di questo lavoro possedeva, e che è del tipo che si scopre soltanto quando un trasferimento fallisce: un esemplare nato dall'uovo misterioso della seconda generazione non passa alla banca se conosce una certa mossa di contatto, per una svista che la fonte spiega con l'assenza di quell'uovo nella versione giapponese su cui lo strumento fu tarato. Appartiene alla stessa famiglia del divieto sulle macchine nascoste, cioè una mossa che blocca un passaggio, e va nella medesima lista di controllo. La stessa pagina fissa inoltre i valori individuali di quell'uovo in due soli insiemi possibili, il che produce un caso raro in cui la fedeltà su quel campo è decidibile in seconda generazione, contro la regola generale che la dichiara indecidibile per difetto di informazione: qui l'informazione esiste perché il gioco non li estrae ma li sceglie fra due.

Due vincoli di ottenibilità toccano infine il profilo di collezione, e vanno conosciuti prima di sceglierlo. Il primo nasce dal fatto che nella seconda generazione un unico insieme di valori determina insieme cromaticità, sesso e forma di Unown, da cui seguono per pura aritmetica due classi che cromatiche non possono essere, cioè le femmine delle specie con rapporto fra i sessi di una a sette e tutte le forme di Unown tranne due lettere; e poiché questo lavoro ha un asse delle forme che conta ventotto voci per quella specie e un asse del sesso che ne conta centodue, il vincolo li attraversa entrambi. Il secondo corregge una simmetria che sarebbe naturale assumere fra i due titoli per console fissa, perché gli esemplari oscuri sono bloccati contro la cromaticità nel secondo e non nel primo. Nella medesima pagina si registrano infine due porte che il deposito stesso apre e che questo lavoro non aveva contato, cioè due mitici consegnati in forma cromatica come premio per il completamento di Pokedex regionali, il che lega un premio al completamento che costituisce l'obiettivo dichiarato di questo lavoro e va verificato contro le condizioni correnti del servizio prima di essere contato.

\section{Un cluster inaccoppiabile per costruzione}

La chiusura ha fatto emergere un difetto nello strumento che porta lo stato di lettura dentro il registro delle fonti, ed è della classe che non produce alcun errore visibile. La prima cella della tabella di stato ammette più nomi di cluster separati da virgola, e il programma la spezzava sempre su quel segno; ma il nome di uno dei cluster contiene a sua volta una virgola, perciò quel cluster non poteva essere accoppiato in nessun modo e sarebbe rimasto dichiarato non letto per sempre, qualunque cosa si fosse scritta nella tabella. Il difetto non si manifesta come un errore ma come una riga che resta ostinatamente nello stato di partenza, che è esattamente il modo in cui un controllo automatico perde la propria utilità mentre continua a funzionare.

La correzione conserva la convenzione invece di cambiarla: la cella si confronta prima per intero con i nomi noti del censimento, e soltanto se non corrisponde a nessuno la si spezza sulla virgola. Chi non passa l'elenco dei nomi noti ottiene il comportamento precedente, cosicché la modifica non alteri il significato delle tabelle già scritte. È la stessa forma di soluzione già adottata altrove in questo lavoro quando una convenzione sintattica entra in conflitto con un dato che la contiene, cioè riconoscere il caso intero prima di applicare la regola generale.
